\section{Keep one cumulative program and prepare separate observations}
Use \texttt{packages/026-chacha20/starter/main.asm} as the cumulative working source. Each authored step replaces instructions only between its named
BEGIN and END comments. Keep labels, caller controls, saves, restores and supplied
status checks. A missing region takes a development exit with status three. A
completed kernel returns zero; the wrapper also uses one and two for its declared
count and counter errors. Development output may be partly changed and must never
be used as a valid keystream or ciphertext.

The caller's mode selects zero for the quarter-round, one for the scalar block,
two for MSA and three for the wrapper. Its \texttt{length} is a full doubleword;
\texttt{in\_place} selects whether the wrapper's destination aliases its source.
The default caller demonstrates separated quarter-round words. It records the
returned status after the \texttt{main\_done} checkpoint and exits through the
ordinary simulator service.

Each recipe specifies exact data, a target and an observation boundary for the
program through that edit. For example, prepare the second quarter-round group
in a separate file using the helper below. Its source argument is your current
working copy, including your completed regions.
\begin{verbatim}
python3 packages/026-chacha20/prepare.py s05-1 \
  --source my-chacha.asm --output observe-s05.asm
\end{verbatim}
Inside an extracted student distribution, run its top-level preparation helper
with your working source explicitly selected. The helper replaces only fixture
arenas and caller controls. It refuses an
existing output file and leaves the authored code intact. Keep the cumulative
source as your editable program; do not use a prepared observation as the next
cumulative source, because its controls are deliberately recipe-specific.

Pause at the recipe's checkpoint before continuing to the temporary or normal
return. Inspect registers inside the frame when requested: their caller values
will be restored on return. The independent construction checks verify both this
inside view and the returned state, using the caller and separate offset arguments.
The strict register mapping describes this worked construction. Final free-writing
checks concern the completed interfaces and allow equivalent internal choices.
